% Independent sidebar overlay; never put answers in the third table cell.
% Oversized or end-of-body pending notes fall back to breakable, full-width
% detailed pages. Their source page keeps a clickable pointer.
\makeatletter
\zref@addprops{savepos}{abspage,page}
\newcommand{\SourceAnchor}[1]{\nobreak\zsavepos{source-#1}}
\def\ReaderNoteList{}
\NewDocumentCommand{\RegisterQA}{m+m+m}{%
  \ReaderRegisterNote{#1}{{\color{ReaderQuestion}\textbf{问题：}#2}\par\smallskip{\color{ReaderAnswer}\textbf{简答：}#3}}}
\NewDocumentCommand{\RegisterAnswer}{m+m}{\ReaderRegisterNote{#1}{{\color{ReaderAnswer}#2}}}
\NewDocumentCommand{\RegisterTutorialPointer}{m+m+m}{%
  \ReaderRegisterNote{#1}{\hyperlink{#2}{\textcolor{ReaderQuestion}{#3}}}}
\long\def\ReaderRegisterNote#1#2{%
  \ifcsname reader-note-content-#1\endcsname
    \PackageError{reader}{Duplicate note ID #1}{Use a unique stable annotation ID.}%
  \fi
  \ifx\ReaderNoteList\@empty\xdef\ReaderNoteList{#1}\else\xdef\ReaderNoteList{\ReaderNoteList,#1}\fi
  \expandafter\gdef\csname reader-note-content-#1\endcsname{#2}%
  \expandafter\gdef\csname reader-note-placed-#1\endcsname{0}%
  \expandafter\newbox\csname reader-note-box-#1\endcsname
  \expandafter\newbox\csname reader-pointer-box-#1\endcsname
  \global\setbox\csname reader-note-box-#1\endcsname=\vbox{%
    \color@begingroup
    \hsize=\dimexpr\NotesWidth-12pt\relax\linewidth=\hsize\columnwidth=\hsize
    \RaggedRight\fontsize{8.6}{12}\selectfont\parindent=0pt\parskip=0pt
    #2\par\color@endgroup}%
  \global\setbox\csname reader-pointer-box-#1\endcsname=\vbox{%
    \color@begingroup
    \hsize=\dimexpr\NotesWidth-12pt\relax\linewidth=\hsize\columnwidth=\hsize
    \RaggedRight\fontsize{7.5}{10}\selectfont\parindent=0pt\parskip=0pt
    \hyperlink{note-#1}{\textcolor{ReaderQuestion}{#1：解答见后续右栏或详细解释页。}}\par\color@endgroup}%
}
\newdimen\ReaderSidebarLeft
\newdimen\ReaderSidebarTop
\newdimen\ReaderSidebarBottom
\newdimen\ReaderSidebarNext
\newdimen\ReaderSourceY
\newdimen\ReaderDesiredY
\newdimen\ReaderNoteHeight
\newbox\ReaderOverlayBox
\newif\ifReaderPending
\newcommand{\ReaderDrawConnector}[2]{%
  \coordinate (reader-src) at ([xshift=\zposx{source-#1}sp,yshift=\zposy{source-#1}sp]current page.south west);
  \coordinate (reader-drop) at ([yshift=-3pt]reader-src);
  \coordinate (reader-gutter) at ([xshift=-3pt]#2.west);
  \draw[ReaderRule,line width=.5pt] (reader-src)--(reader-drop)--(reader-gutter |- reader-drop)--(reader-gutter |- #2.west)--(#2.west);
}
\newcommand{\ReaderTryNote}[1]{%
  \edef\ReaderSourcePage{\zref@extractdefault{source-#1}{abspage}{0}}%
  \ifnum\csname reader-note-placed-#1\endcsname=0\relax
    \ifnum\ReaderSourcePage>0\relax
      \ifnum\ReaderSourcePage>\ReadonlyShipoutCounter\relax\else
        \ReaderNoteHeight=\dimexpr\ht\csname reader-note-box-#1\endcsname+\dp\csname reader-note-box-#1\endcsname+12pt\relax
        \ReaderDesiredY=\ReaderSidebarNext
        \ifnum\ReaderSourcePage=\ReadonlyShipoutCounter\relax
          \ReaderSourceY=\dimexpr\paperheight-\zposy{source-#1}sp-12pt\relax
          \ifdim\ReaderSourceY>\ReaderDesiredY\ReaderDesiredY=\ReaderSourceY\fi
          % Moving the complete box upward is permitted; body remains fixed.
          \ifdim\dimexpr\ReaderDesiredY+\ReaderNoteHeight\relax>\ReaderSidebarBottom
            \ReaderDesiredY=\dimexpr\ReaderSidebarBottom-\ReaderNoteHeight\relax
            \ifdim\ReaderDesiredY<\ReaderSidebarNext\ReaderDesiredY=\ReaderSidebarNext\fi
          \fi
        \fi
        \ifdim\dimexpr\ReaderDesiredY+\ReaderNoteHeight\relax>\ReaderSidebarBottom
          % Leave the note queued; a later body page or the appendix will hold it.
        \else
          \global\setbox\ReaderOverlayBox=\copy\csname reader-note-box-#1\endcsname
          \node[anchor=north west,draw=ReaderRule,line width=.5pt,fill=white,
            inner xsep=4pt,inner ysep=6pt,outer sep=0pt,align=left]
            (reader-note-#1) at ([xshift=\ReaderSidebarLeft,yshift=-\ReaderDesiredY]current page.north west)
            {\hypertarget{note-#1}{}\box\ReaderOverlayBox};
          \expandafter\gdef\csname reader-note-placed-#1\endcsname{1}%
          \ifnum\ReaderSourcePage=\ReadonlyShipoutCounter\relax\ReaderDrawConnector{#1}{reader-note-#1}\fi
          \immediate\write\ReaderNoteLog{#1|\the\ReadonlyShipoutCounter|sidebar|\strip@pt\ReaderDesiredY|\strip@pt\ReaderNoteHeight}%
          \ReaderSidebarNext=\dimexpr\ReaderDesiredY+\ReaderNoteHeight+12pt\relax
        \fi
      \fi
    \fi
  \fi
  \ifnum\ReaderSourcePage=\ReadonlyShipoutCounter\relax
    \ifnum\csname reader-note-placed-#1\endcsname=0\relax
      % Pending pointers use the same allocation cursor as QA frames. When no
      % room remains, the highlighted source itself is still a direct link.
      \ReaderSourceY=\dimexpr\paperheight-\zposy{source-#1}sp\relax
      \ReaderNoteHeight=\dimexpr\ht\csname reader-pointer-box-#1\endcsname+\dp\csname reader-pointer-box-#1\endcsname+4pt\relax
      \ReaderDesiredY=\ReaderSidebarNext
      \ifdim\ReaderSourceY>\ReaderDesiredY\ReaderDesiredY=\ReaderSourceY\fi
      \ifdim\dimexpr\ReaderDesiredY+\ReaderNoteHeight\relax>\ReaderSidebarBottom
        \ReaderDesiredY=\dimexpr\ReaderSidebarBottom-\ReaderNoteHeight\relax
        \ifdim\ReaderDesiredY<\ReaderSidebarNext\ReaderDesiredY=\ReaderSidebarNext\fi
      \fi
      \ifdim\dimexpr\ReaderDesiredY+\ReaderNoteHeight\relax>\ReaderSidebarBottom\else
        \global\setbox\ReaderOverlayBox=\copy\csname reader-pointer-box-#1\endcsname
        \node[anchor=north west,inner sep=0pt,outer sep=0pt,align=left]
          (reader-pointer-#1) at ([xshift=\ReaderSidebarLeft,yshift=-\ReaderDesiredY]current page.north west)
          {\box\ReaderOverlayBox};
        \ReaderDrawConnector{#1}{reader-pointer-#1}%
        \immediate\write\ReaderNoteLog{#1|\the\ReadonlyShipoutCounter|pointer|\strip@pt\ReaderDesiredY|\strip@pt\ReaderNoteHeight}%
        \ReaderSidebarNext=\dimexpr\ReaderDesiredY+\ReaderNoteHeight+8pt\relax
      \fi
    \fi
  \fi
}
\AddToHook{shipout/foreground}{%
  \ifReaderDetailActive\else
    \begingroup
      \ReaderSidebarLeft=\dimexpr1in+\hoffset+\oddsidemargin+5\tabcolsep+\EnglishWidth+\ChineseWidth+\arrayrulewidth\relax
      \ReaderSidebarTop=\dimexpr1in+\voffset+\topmargin+\headheight+\headsep\relax
      \ReaderSidebarBottom=\dimexpr\ReaderSidebarTop+\textheight-4pt\relax
      \ReaderSidebarNext=\ReaderSidebarTop
      \begin{tikzpicture}[remember picture,overlay]
        \@for\ReaderID:=\ReaderNoteList\do{\expandafter\ReaderTryNote\expandafter{\ReaderID}}%
      \end{tikzpicture}
    \endgroup
  \fi
}
\newcommand{\ReaderFlushNotes}{%
  \clearpage
  \ReaderPendingfalse
  \@for\ReaderID:=\ReaderNoteList\do{%
    \ifnum\csname reader-note-placed-\ReaderID\endcsname=0\relax\ReaderPendingtrue\fi}%
  \ifReaderPending
    \ReaderDetailedPage{reader-pending-notes}{阅读备注：补充解答}{%
      \@for\ReaderID:=\ReaderNoteList\do{%
        \ifnum\csname reader-note-placed-\ReaderID\endcsname=0\relax
          \par\medskip\hypertarget{note-\ReaderID}{}%
          {\bfseries\ReaderID}\quad
          \edef\ReaderSourcePage{\zref@extractdefault{source-\ReaderID}{page}{?}}%
          \edef\ReaderSourceAbsolute{\zref@extractdefault{source-\ReaderID}{abspage}{0}}%
          \ifnum\ReaderSourceAbsolute=0\relax
            \PackageWarningNoLine{reader}{Unresolved source anchor \ReaderID. Rerun LaTeX; if persistent, add SourceAnchor for this ID}%
          \fi
          对应正文第\ReaderSourcePage 页\par\smallskip
          \csname reader-note-content-\ReaderID\endcsname\par\medskip
          \expandafter\gdef\csname reader-note-placed-\ReaderID\endcsname{1}%
          \immediate\write\ReaderNoteLog{\ReaderID|appendix|detail||}%
        \fi}%
    }%
  \fi
}
% This hook is a final safety net. Calling ReaderFlushNotes explicitly after
% the final table is clearer and allows body content to follow if desired.
\AtEndDocument{\ReaderFlushNotes\immediate\closeout\ReaderNoteLog\immediate\closeout\ReaderPageLog}
\makeatother
